% SEGMENT OLP-0429-S01
% Part: normal-modal-logic
% Chapter: axioms-systems
% Section: normal-logics

\documentclass[../../../include/open-logic-section]{subfiles}

\begin{document}

\olfileid{nml}{prf}{nor}

\olsection{இயல்பான வாய்ப்புநிலைத் தருக்கங்கள்}

வாய்ப்புநிலை !!{formula}s இன் ஒவ்வொரு கணத்தையும், ஓர் அடிகோள் கணத்திலிருந்து
வருவிக்கத்தக்க !!{formula}s இன் கணமாக எளிதில் சிறப்பிக்க முடியாது.
வாய்ப்புநிலைத் தருக்கங்கள் நன்கு நடந்துகொள்பவையாக இருக்க வேண்டும் என்று
விரும்புகிறோம். முதலில், செவ்வியல் கூற்றுத் தருக்கத்தில் நாம் !!{derive}
செய்யக்கூடிய அனைத்தும் இன்னும் !!{derivable} ஆக இருக்க வேண்டும்; இப்போது
!!{formula}s இல்
\iftag{prvBox}{$\Box$\iftag{prvDiamond}{
    மற்றும்~}{}}{}\iftag{prvDiamond}{$\Diamond$}{} உம் வரலாம் என்பதை
நிச்சயமாகக் கணக்கில் கொள்ள வேண்டும். இதற்காக, ஒரு வாய்ப்புநிலைத் தருக்கம்
எல்லா மெய்ம நிகழ்வுகளையும் கொண்டிருக்க வேண்டும் என்றும், மோடஸ்
போனென்ஸின் கீழ் மூடப்பட்டிருக்க வேண்டும் என்றும் கோருகிறோம்.

\begin{defn}
  ஒரு \emph{வாய்ப்புநிலைத் தருக்கம்} என்பது பின்வருமாறு அமையும்
  வாய்ப்புநிலை !!{formula}s இன் கணம்~$\Sigma$:
  \begin{enumerate}
  \item அது எல்லா மெய்மங்களையும் கொண்டுள்ளது; மேலும்
  \item அது பதிலீட்டின் கீழ் மூடப்பட்டுள்ளது; அதாவது $!A \in \Sigma$,
    மேலும் $!D_1$, \dots,~$!D_n$ ஆகியவை !!{formula}s எனில்,
    \[
    \SSubst{!A}{\subst{!D_1}{p_1}, \dots, \subst{!D_n}{p_n}} \in \Sigma,
    \]
  \item அது \emph{மோடஸ் போனென்ஸின்} கீழ் மூடப்பட்டுள்ளது; அதாவது $!A$,
    $!A \lif !B \in \Sigma$ எனில், $!B \in \Sigma$.
  \end{enumerate}
\end{defn}

வாய்ப்புநிலைத் தருக்கங்களுக்குத் தொடர்புசார் பொருண்மையியலைப் பயன்படுத்த,
எல்லா வாய்ப்புநிலை மாதிரிகளிலும் செல்லுபடியாகும் எல்லா !!{formula}s உம்
சேர்க்கப்பட வேண்டும் என்றும் கோர வேண்டும். \Ax{K}\iftag{prvDiamond}{
மற்றும்~\Dual{}}{} இன் எல்லா நிகழ்வுகளும் !!{derivable} ஆகவும், ஒரு
!!a{formula}~$!A$ !!{derivable} ஆகும்போதெல்லாம்~$\Box !A$ உம் அவ்வாறு
ஆகவும் இருந்தாலே இக்கோரிக்கை நிறைவுறுகிறது என்று தெரியவருகிறது. இந்நிபந்தனைகளை
நிறைவுசெய்யும் ஒரு வாய்ப்புநிலைத் தருக்கம் \emph{இயல்பானது} எனப்படுகிறது.
(இயல்பற்ற வாய்ப்புநிலைத் தருக்கங்களும் உண்டு; ஆனால் வழக்கமான தொடர்புசார்
மாதிரிகள் அவற்றுக்குப் போதுமானவை அல்ல.)

\begin{defn}
  ஒரு வாய்ப்புநிலைத் தருக்கம்~$\Sigma$ பின்வருவனவற்றைக் கொண்டிருந்தால்
  \emph{இயல்பானது}:
  \begin{align*}
    \tag{\Ax{K}} & \Box(p \lif q) \lif (\Box p \lif \Box q),
    \iftag{prvDiamond}{\\
      \tag{\Dual} & \Diamond p \liff \lnot\Box\lnot p}{}
  \end{align*}
  மேலும் அது \emph{இன்றியமையாக்கலின்} கீழ் மூடப்பட்டிருக்க வேண்டும்;
  அதாவது $!A \in \Sigma$ எனில், $\Box !A \in \Sigma$.
\end{defn}

மெய்ம உட்கிடைப்பு \emph{ஓர் உலகில் மெய்மையைப்} பேணும் அளவுக்கு
“நுணுக்கமானது”; ஆனால் விதி~\Nec{} \emph{ஒரு மாதிரியில் மெய்மையை} மட்டுமே
பேணுகிறது (எனவே ஒரு சட்டகத்தில் அல்லது சட்டகங்களின் ஒரு வகுப்பில்
செல்லுபடித்தன்மையையும் பேணுகிறது) என்பதைக் கவனிக்க.

\begin{prop}\ollabel{prop:rk}
  ஒவ்வொரு இயல்பான வாய்ப்புநிலைத் தருக்கமும் பின்வரும் விதி~\RK{} இன் கீழ்
  மூடப்பட்டுள்ளது:
  \begin{prooftree}
    \AxiomC{$!A_1 \lif (!A_2 \lif \cdots (!A_{n-1} \lif !A_n)\cdots)$}
    \RightLabel{\RK}
    \UnaryInfC{$\Box!A_1 \lif (\Box!A_2 \lif \cdots (\Box!A_{n-1}
      \lif \Box!A_n)\cdots).$}
  \end{prooftree}
\end{prop}

\begin{proof}
  $n$ மீதான தொகுத்தறிதலால்: $n = 1$ எனில், இவ்விதி~\Nec{} தான்;
  ஒவ்வொரு இயல்பான வாய்ப்புநிலைத் தருக்கமும்~\Nec{} இன் கீழ்
  மூடப்பட்டுள்ளது.

  இப்போது முடிவு $n-1$ க்கு உண்மை எனக் கொண்டு, அது~$n$ க்கும் உண்மை என்று
  காட்டுவோம்.

  பின்வருவதை எடுத்துக்கொள்க:
  \begin{align*}
  & !A_1 \lif (!A_2 \lif \cdots (!A_{n-1} \lif !A_n)\cdots) \in \Sigma
  \intertext{தொகுத்தறிதல் கருதுகோளின்படி, நமக்குப் பின்வருவது உள்ளது:}
  & \Box!A_1 \lif (\Box!A_2 \lif \cdots \Box(!A_{n-1} \lif !A_n)\cdots)
  \in \Sigma
  \intertext{$\Sigma$ ஓர் இயல்பான வாய்ப்புநிலைத் தருக்கம் என்பதால்,
    $\Ax{K}$ இன் எல்லா நிகழ்வுகளையும், குறிப்பாகப் பின்வருவதையும், கொண்டுள்ளது:}
  & \Box(!A_{n-1} \lif !A_n) \lif (\Box!A_{n-1} \lif \Box!A_n) \in \Sigma
  \intertext{மோடஸ் போனென்ஸையும் பொருத்தமான மெய்ம நிகழ்வுகளையும் பயன்படுத்தி,
    பின்வருவதைப் பெறுகிறோம்:}
  & \Box!A_1 \lif (\Box!A_2 \lif \cdots (\Box!A_{n-1}
  \lif \Box!A_n)\cdots) \in \Sigma.
  \end{align*}
\end{proof}

\begin{prop}\ollabel{prop:notDiamondBot}
  ஒவ்வொரு இயல்பான வாய்ப்புநிலைத் தருக்கம்~$\Sigma$ உம்
  $\lnot\Diamond\lfalse$ ஐக் கொண்டுள்ளது.
\end{prop}

\begin{prob}
  \olref[nml][prf][nor]{prop:notDiamondBot} ஐ நிறுவுக.
\end{prob}

% SOURCE CORRECTION TA-NML-015: the proof intersects normal modal logics, so the claimed least logic must also be normal.
\begin{prop}
  $!A_1$, \dots,~$!A_n$ ஆகியவை !!{formula}s ஆக இருக்கட்டும். $!A_1$,
  \dots,~$!A_n$ ஆகியவற்றின் எல்லா நிகழ்வுகளையும் கொண்ட மிகச்சிறிய இயல்பான
  வாய்ப்புநிலைத் தருக்கம்~$\Sigma$ ஒன்று உள்ளது.
\end{prop}

\begin{proof}
  $!A_1$, \dots,~$!A_n$ கொடுக்கப்பட்டால், $!A_1$, \dots,~$!A_n$
  ஆகியவற்றின் எல்லா நிகழ்வுகளையும் கொண்ட எல்லா இயல்பான வாய்ப்புநிலைத்
  தருக்கங்களின் வெட்டாக~$\Sigma$ ஐ வரையறுக்க. எல்லா !!{formula}s இன்
  கணமான~$\Frm[L]$ அத்தகைய ஒரு
  வாய்ப்புநிலைத் தருக்கம் என்பதால், இந்த வெட்டு காலியல்ல.
\end{proof}

\begin{defn}
$!A_1$, \dots,~$!A_n$ ஆகியவற்றைக் கொண்ட மிகச்சிறிய இயல்பான வாய்ப்புநிலைத்
தருக்கம் ஒரு \emph{வாய்ப்புநிலை அமைப்பு} எனப்படுகிறது; அது
$\Log{K} !A_1 \dots !A_n$ என்று குறிக்கப்படுகிறது. மிகச்சிறிய இயல்பான
வாய்ப்புநிலைத் தருக்கம்~\Log{K} என்று குறிக்கப்படுகிறது.
\end{defn}

\end{document}
